Recent Large Language Models (LLMs) have been training on a diverse range of data sources, which presumable include works of 
mathematical science and logical theories such as Natural Deduction and Category Theory. It has already been demonstrated
that these models can understand symbolic notation \autocite{vangassen2026}.

The inclusion of such material in their training datasets raises several 
questions regarding the nature of their reasoning capabilities and the potential implications for use cases in the 
field of artificial intelligence.

This paper is a empirical investigation into how well selected LLMs can understand and apply symbolic notation and concepts from mathematical science
and logical theories. The experiment is run of a set of 42 selected LLMs ranging from different parameter sizes, vendors and architectures. 
Each LLM is presented with a prompt formulated in symbolic notation; a story context and few rules; and a question regarding 
the presented context that requires an application of the presented rules. 

\subsection{Research Questions}

This paper tries to adress the following research questions:
\begin{itemize}
    \item How well do LLMs understand and apply symbolic notation from mathematical science and logical theories?
    \item How well can LLMs output results in symbolic notation instead of prose language?
\end{itemize}

The ability for both ingesting and outputting results in symbolic notation is important for the development of 
agentic systems \autocite{ross2026aisp} which require stricter formalism and which want to explicitely keep prose language out of their processes.

To address these questions, we designed an empirical experiment using a standardized Python testbed and a set of LLMs from diverse vendors and architectures.